\section{Methods}
\label{sec:methods}

\subsection*{Definitions and classification}

We calculate $R_C=C_{\mathrm{fast}}/W$ from a declared usable fast-tier
capacity and active working-set size. We calculate
$R_B=T_{\mathrm{comp}}/T_{\mathrm{transfer}}$, with compute and transfer timed
separately. Classification gives capacity precedence: $R_C<0.5$ is
capacity-limited; otherwise $R_B<1$ is I/O-limited; all remaining points are
coordination-dominated. This rule is implemented in the released Python code.

\subsection*{Pointer-chain capacity probe}

The CPU probe allocates a NumPy array and constructs a seeded random
permutation. A C11 kernel compiled with \texttt{-O3} traverses it as a dependent
pointer chain; each loaded index determines the address of the next load,
limiting vectorisation, instruction-level parallelism, and prefetching. We
measured working sets of 0.5, 1, 2, 4, 8, 16, 24, 32, 48, 64, 96, 128, 192,
and 256~MiB. Each point used seven timed trials after warm-up. The output stores
working-set size, $R_C$, mean nanoseconds per dependent load, standard
deviation, and trial count. The fast-tier capacity was read from Linux sysfs
and was 48~MiB. The host exposed three logical CPUs of an Intel Xeon Platinum
8370C under KVM; the summary records the kernel, Python, and NumPy versions.
Timing used \texttt{clock\_gettime(CLOCK\_MONOTONIC\_RAW)} inside the compiled
loop. No outlier was removed and CPU affinity or frequency was not controlled.

\subsection*{Layered matrix and copy probe}

The second probe allocates 12 square, float32 layer matrices with $d=512$ and
batch size 8. A selected fraction remains resident; non-resident matrices are
copied before multiplication. NumPy's wall-clock path was used on the CPU, with
five windows per point. Matrix entries are deterministically seeded and scaled
by $1/\sqrt{d}$ to prevent numerical overflow during repeated multiplication.
The script also supports CUDA-event timing and XLA synchronisation, but no GPU
or TPU result is reported here.

For the reported capacity ratio, we averaged total latency at the three sampled
points below $R_C=0.5$ and divided by latency at $R_C=1$. This endpoint summary
is descriptive; no hypothesis test or population-level confidence interval is
claimed. The pointer-chain table reports mean and population standard deviation
across the seven repeated trials exactly as stored in the JSON summary.

\subsection*{Software and provenance}

The compiled pointer-chain run used Python 3.12.13, NumPy 2.5.2, and GCC 13.3;
the earlier layered-matrix harness check used Python 3.11.15, but that summary
did not capture the NumPy version. Seeds, grids, timing code, raw records, summaries, and
commands are version controlled. The analytical regime
map reads threshold constants from the same module as the classifier. The
simulated backend is provided solely for software testing and qualitative
exploration; its outputs are excluded from the reported measurements.

\subsection*{Statistics and reporting}

This proof-of-concept study did not use random assignment, blinding, or a
sample-size calculation. Trial counts were fixed before summarisation. We
report all measured grid points in the machine-readable data, selected points
in Table~\ref{tab:cpu_probe}, arithmetic means, and one standard deviation. No
claim of statistical significance is made. Hardware-platform generalisation
requires independent machines and is left for future work.
